\documentclass[11pt]{article}

\usepackage[margin=1in]{geometry}
\usepackage{amsmath,amssymb}
\usepackage{booktabs}
\usepackage{microtype}
\usepackage[hidelinks]{hyperref}
\usepackage{enumitem}
\setlist{nosep,leftmargin=1.5em}

\newcommand{\dk}{d/d k}
\newcommand{\dt}{\partial_t}
\newcommand{\Lt}{\tilde\Lambda}
\newcommand{\Gt}{\tilde G}
\newcommand{\lfour}{\lambda_4}
\newcommand{\gp}{\bar{g}}

\title{A Declared-Truncation Reconstruction of a 41-Coupling
Gravity--Standard-Model FRG System: Verified Blocks, Joint Fixed Point, and
Stability Spectrum}

\author{C.~M.~Baird \and Zora}
\date{\today}

\begin{document}
\maketitle

\begin{abstract}
We reconstruct, in a fully declared and auditable manner, the running-coupling
system of a four-dimensional functional-renormalization-group (FRG) truncation
that couples quantum gravity to Standard-Model matter and two real gauge-singlet
scalars at mass dimension $d\le 6$.  The original archival claim of a
$41$-running-coupling ultraviolet completion is not accompanied by an operator
list, a beta-functional system, a fixed-point vector, or a stability matrix;
rather than fabricating the missing input, we declare a $41$-coordinate basis,
record a regulator, gauge, and projection convention for every block, and
derive only those flows that trace to a written source or a documented
projection.  Thirty-three of the forty-one beta functions are thereby supplied
and verified: the Einstein--Hilbert--matter block of Don\`a, Eichhorn and
Percacci; the nonminimal scalar--curvature system of Narain and Percacci; the
scalar local-potential block at quartic and sixth order; and the
curvature-squared block of Codello, Percacci and Rahmede.  The joint fixed
point of the derived flow is the DEP Standard-Model fixed point
$(\lambda_N,g_N)=(-4.7975,1.7621)$ with Gaussian matter.  The $33$-dimensional
quantum-plus-canonical stability spectrum contains 8 relevant, 12 marginal,
and 13 irrelevant directions; the two gravity exponents reproduce the
published values $(3.961,1.644)$.  We also record a negative result: in the
one-loop curvature-squared block the scaffold coordinate of the Weyl-squared
coupling runs monotonically, so that block admits no fixed point in that
parametrization.  Six derivative and gauge-portal operators remain undriven
and are explicitly declared pending.
\end{abstract}

\section{Introduction}
\label{sec:intro}

The asymptotic-safety scenario for quantum gravity
\cite{reuter,litim06,percacciBook} is one of the few quantum-gravity proposals
that can be addressed quantitatively.  Within the functional renormalization
group, the central object is the effective average action $\Gamma_k$, whose
scale dependence is governed by the exact Wetterich equation
\cite{wetterich,berges},
\begin{equation}
  \label{eq:wetterich}
  \dt\Gamma_k[\Phi] = \frac12\,\mathrm{STr}\!\left[
    \left(\Gamma_k^{(2)}[\Phi]+R_k\right)^{-1}\dt R_k\right],
\end{equation}
where $\Gamma_k^{(2)}$ is the second functional derivative (Hessian), $R_k$ a
momentum cutoff (regulator), and $t=\ln(k/k_0)$ the logarithmic scale
parameter.  In practice the equation is solved in truncated theory spaces, and
every numerical claim carries a hidden load of choices: the operator basis,
the regulator, the gauge fixing and ghost sector, and the projection
background.

This paper is the scientific write-up of a reconstruction exercise that began
from a corpus claiming a $41$-running-coupling, dimension-$d\le 6$ truncation
coupling gravity to the Standard Model plus two scalar singlets.  The
claim was not accompanied by the mathematical objects needed to reproduce it:
no complete operator list, no beta functions, no fixed-point vector, and no
$41\times 41$ stability matrix.  The reconstruction therefore proceeds under a
strict \emph{no-fabrication} protocol: every beta function must trace to a
written source equation (with the paper, equation, gauge, and regulator
recorded) or to an explicitly declared projection of a documented trace
functional.  Coordinates without a traced flow are marked \texttt{NaN} (not
set to zero), and the framework refuses to present the result as the
historical $41$-coupling theory.  Instead it delivers a smaller, fully
declared truncation whose content is \emph{verified}, together with an honest
inventory of what remains undriven.

The main results of this paper are:
\begin{enumerate}
  \item a declared $41$-coordinate basis with recorded conventions
        (Sec.~\ref{sec:framework}, App.~\ref{app:operators});
  \item five verified building blocks, each checked against its published
        source (Sec.~\ref{sec:blocks});
  \item the joint fixed point of the derived $33$-coordinate flow and its
        stability spectrum, with 8 relevant, 12 marginal, and 13 irrelevant
        directions (Sec.~\ref{sec:joint});
  \item a recorded negative result in the curvature-squared sector
        (Sec.~\ref{sec:r2}, Sec.~\ref{sec:discussion});
  \item an explicit list of the six pending operators that remain undriven
        (Sec.~\ref{sec:coverage}).
\end{enumerate}

Throughout, the numbers quoted are produced by the accompanying Python
implementation, which is exercised by a unit-test suite; every table in this
paper corresponds to a test or a fixed-point verification in that suite.

\section{Framework, conventions, and controls}
\label{sec:framework}

\subsection{Fields, action, and the Wetterich equation}
\label{sec:fields}

We work in Euclidean signature in $d=4$ dimensions with a background-field
formulation.  The field content is the metric $g_{\mu\nu}$, the Standard-Model
fields, and two real gauge-singlet scalars $\Phi_c$ and $E$.  The flow
equation is the Wetterich equation \eqref{eq:wetterich}; the Hessian is
constructed covariantly, background-covariant gauge fixing with
Faddeev--Popov ghosts is used, and a specific regulator and projection
background must be recorded with every numerical result.  The implementation
adopts the optimized (Litim) shape \cite{litim01},
\begin{equation}
  \label{eq:litim}
  R_k(q^2)=\left(k^2-q^2\right)\theta\!\left(k^2-q^2\right),
\end{equation}
except in the gravity--matter block of Sec.~\ref{sec:dep}, which uses the
type-II cutoff of its source paper; the two regulator families are never mixed
in one solve.

\subsection{The declared 41-coordinate basis}
\label{sec:basis}

The coordinates follow \texttt{model.NAMES} in order:
\begin{enumerate}[label=\arabic*.]
  \item $g_N$: dimensionless Newton coupling (Einstein--Hilbert, $R$).
  \item $\lambda_N$: dimensionless cosmological term.
  \item $a_{R^2}$, $a_{C^2}$: $R^2$ and Weyl-squared coordinates.
  \item--7.\quad $g_1,g_2,g_3$: Standard-Model gauge couplings.
  \item--10.\quad $y_t,y_b,y_\tau$: third-generation Yukawas.
  \item--16.\quad $\lambda_H,m_H^2$; $\lambda_\Phi,m_\Phi^2$; $\lambda_E,m_E^2$:
        quartics and masses of the Higgs and the two singlets.
  \item--19.\quad $\kappa_{\Phi E}$, $\kappa_{H\Phi}$, $\kappa_{HE}$:
        quartic portals $\Phi_c^2E^2$, $|H|^2\Phi_c^2$, $|H|^2E^2$.
  \item--22.\quad $\xi_H,\xi_\Phi,\xi_E$: nonminimal couplings
        $R|H|^2$, $R\Phi_c^2$, $R E^2$.
  \item--25.\quad $Z_h,Z_\phi,Z_E$: wave-function renormalizations.
  \item--35.\quad ten scalar dimension-six potentials: $|H|^6$, $\Phi^6$,
        $E^6$, $|H|^4\Phi^2$, $|H|^4E^2$, $|H|^2\Phi^4$, $|H|^2E^4$,
        $\Phi^4E^2$, $\Phi^2E^4$, $|H|^2\Phi^2E^2$ (implementation names
        $c_{H6}\ldots c_{H2\Phi2E2}$).
  \item--38.\quad derivative dimension-six operators: $|H|^2\Box|H|^2$,
        $\Phi^2\Box\Phi^2$, $E^2\Box E^2$ ($c_{H\Box},c_{\Phi\Box},c_{E\Box}$).
  \item--41.\quad gauge-strength portals $|H|^2F^2$, $\Phi^2F^2$, $E^2F^2$
        ($c_{HF^2},c_{\Phi F^2},c_{EF^2}$).
\end{enumerate}
The full table with physical definitions is collected in
App.~\ref{app:operators}.

\subsection{Canonical scaling and the no-fabrication control}
\label{sec:canonical}

Every dimensionless coupling receives its canonical contribution.  For a
coupling $g_i$ whose mass dimension is $[g_i]$, the canonical term is
$- [g_i]\, g_i$; the non-vanishing entries are
\begin{equation}
  \label{eq:canonical}
  \beta_{g_N}\supset +2\,g_N,\qquad
  \beta_{\lambda_N}\supset -2\,\lambda_N,\qquad
  \beta_{m_i^2}\supset -2\,m_i^2,\qquad
  \beta_{c_\bullet}\supset +2\,c_\bullet ,
\end{equation}
with $c_\bullet$ any dimension-six coordinate (the projected sextic
couplings carry mass dimension $-2$, since the background field carries mass
dimension $1$).  All other coordinates are dimensionless ($[g_i]=0$).

The no-fabrication control is implemented at the code level: assembling the
full $41$-vector beta function without a complete projected quantum
contribution raises \texttt{IncompleteFlowError}; entries for which no traced
projection exists are \texttt{NaN}, never silently zero.  A valid fixed-point
claim must make all $41$ residuals small, and the stability matrix must be the
Jacobian of the declared flow with a stated critical-exponent sign
convention.

\subsection{Residuals, stability matrices, and critical exponents}
\label{sec:methods}

All stability results in this paper use the numerical central-difference
Jacobian of the declared flow, evaluated at the point of interest with step
$\delta$; the critical exponents are the eigenvalues of the Jacobian with the
convention
\begin{equation}
  \label{eq:theta}
  \theta = -\mathrm{eig}\!\left(\partial \beta_i/\partial g_j\right),
\end{equation}
so that a positive real part means a relevant (UV-attractive) direction.
Fixed-point residuals are quoted as $|\beta_i|$ at the point.

\section{Verified building blocks}
\label{sec:blocks}

\subsection{Einstein--Hilbert gravity coupled to Standard-Model matter
(DEP block)}
\label{sec:dep}

The gravity--matter block follows Don\`a, Eichhorn and Percacci
\cite{dep14} (arXiv:1311.2898).  Its convention label is
\texttt{DEP-TYPEII}: type-II cutoff with shape parameters $\alpha=\rho=1$, the
table ghost factor $2\,2^{[d/2]}$, and minimal coupling of matter.  The block
uses the scalar--fermion--gauge content of the Standard Model ($n_s=4$ Higgs
real components, $n_f$ fermions, $n_v$ gauge fields); in this reconstruction
$n_s=4$, $n_f=45/2$ per the standard count, and $n_v=12$, plus the two real
singlets.  The anomalous dimensions of the DEP block are the coupled solution
$\eta_h,\eta_\phi,\eta_E,\eta_V$ of the source equations.

In the scaffold coordinates of this project, $g_N=\Gt$ and
$\lambda_N=2\Lt$, so
\begin{equation}
  \label{eq:scaffold}
  \beta_{g_N}=\beta_{\Gt},\qquad
  \beta_{\lambda_N}=2\,\beta_{\Lt},
\end{equation}
where $(\Lt,\Gt)=(\Lambda k^{-2}, G k^{d-2})$ are the couplings of
Ref.~\cite{dep14}.  The block reproduces the published pure-gravity fixed
point $(\Lt,\Gt)=(0.01046,0.7725)$ and the Standard-Model row
$(\Lt,\Gt)=(-2.399,1.762)$.  In the unimproved (table) treatment used here the
Standard-Model fixed point is
\begin{equation}
  \label{eq:depfp}
  \Lt_*=-2.3987548,\qquad \Gt_*=1.7621271,\qquad
  \theta=(3.9609,\;1.6436),
\end{equation}
with $\eta_h\simeq 0.269$ and $\eta_s\simeq -0.076$ at the fixed point.  The
improved (self-consistent $\eta$) treatment yields
$(\Lt_*,\Gt_*)=(-2.4644,1.7951)$ and $\theta=(3.9988,1.6377)$; the table
convention is used throughout this paper because it is the convention under
which the source's published exponent table is reproduced.  All numbers are
regulator- and gauge-dependent.

\subsection{The Litim Einstein--Hilbert benchmark}
\label{sec:litim}

As a cross-check of the fixed-point machinery, the pure Einstein--Hilbert
truncation with the optimized cutoff \cite{litim01,litim06} is implemented in
its own normalization ($g_{eq}=c_d\,g_N$, $c_4=24\pi$, $\alpha\to\infty$).
The known result is reproduced exactly:
\begin{equation}
  \label{eq:litimfp}
  \lambda_N^*=\frac14,\qquad g_N^*=\frac{1}{1536\pi},\qquad
  \theta=\frac53\pm \frac{i\sqrt{167}}{3}.
\end{equation}
This block is kept deliberately separate from the DEP convention; its only
role is to certify the Jacobian and eigenvalue pipeline.

\subsection{Nonminimal scalar--curvature couplings (Narain--Percacci block)}
\label{sec:np}

The nonminimal couplings $\xi_H,\xi_\Phi,\xi_E$ follow the closed-form
five-coupling scalar--tensor system of Narain and Percacci
\cite{narain10} (arXiv:0911.0386), convention label \texttt{NP-DE-DONDER-TYPEA}:
$d=4$, de-Donder gauge, optimized cutoff, one real scalar.  The couplings are
defined by
\begin{equation}
  \label{eq:npaction}
  V=k^d\sum_{n}\lambda_{2n}\,\tilde\phi^{2n},\qquad
  F=k^{d-2}\sum_{n}\xi_{2n}\,\tilde\phi^{2n},\qquad
  \tilde\phi=\phi\,k^{-(d-2)/2},
\end{equation}
so that in $d=4$, $\lambda_2=m^2/k^2$ and $\lambda_4$ is the dimensionless
quartic coefficient of $\phi^4$.  The system is implicit: the anomalous
dimension $\eta=\dt\xi_0/\xi_0$ is fixed by the $\xi_0$ flow, and
$\dt\xi_2$ appears linearly in the $\lambda_2$, $\xi_2$ and $\lambda_4$
flows.  Writing $\beta_{\xi_0}=X+\eta\,Y$, the implicit condition
$\dt\xi_0=(\eta-2)\xi_0$ is solved algebraically as
\begin{equation}
  \label{eq:solveeta}
  \eta=\frac{X+2\xi_0}{\xi_0-Y},
\end{equation}
with ($\lambda_0,\xi_0$ the dimensionless $R$ and $R\phi^2$ couplings of the
paper)
\begin{align}
  \label{eq:XY}
  X&=-2\xi_0+\frac{1}{384\pi^2}\Bigg[
    25-\frac{4}{1+2\lambda_2}
    -\frac{24\,\xi_2}{(1+2\lambda_2)^2}
    +\frac{8\xi_0\,(7\xi_0-2\lambda_0)}{(\xi_0-\lambda_0)^2}\Bigg],\\
  Y&=\frac{1}{1152\pi^2}\,
    \frac{17\xi_0^2+18\xi_0\lambda_0-15\lambda_0^2}{(\xi_0-\lambda_0)^2},
\end{align}
and
\begin{align}
  \label{eq:betalam0}
  \beta_{\lambda_0}&=-4\lambda_0+\frac{1}{32\pi^2}
    \Bigg[2+\frac{1}{1+2\lambda_2}+\frac{6\lambda_0}{\xi_0-\lambda_0}\Bigg]
    +\frac{\eta}{96\pi^2}\,\frac{5\xi_0-2\lambda_0}{\xi_0-\lambda_0},\\
  \label{eq:betaxi0}
  \beta_{\xi_0}&=-2\xi_0+X+\eta\,Y .
\end{align}
The remaining three flows $\beta_{\lambda_2}$, $\beta_{\xi_2}$, $\beta_{\lfour}$
are closed forms in $(\lambda_0,\xi_0,\lambda_2,\xi_2,\lfour,\eta)$; the
$\xi_2$ flow is solved for the implicit $\dt\xi_2$ term,
\begin{equation}
  \label{eq:xi2implicit}
  \beta_{\xi_2}=\frac{Z+\eta\,W}{1-V/\xi_0},
\end{equation}
and $\beta_{\lambda_2}$, $\beta_{\lfour}$ contain the solved $\dt\xi_2$
combination.  The full forms are listed in
App.~\ref{app:np}.

The two-coupling gravity sector (GMFP) of the paper is reproduced exactly:
\begin{equation}
  \label{eq:gmfp}
  \lambda_0^*=0.0086200674,\qquad \xi_0^*=0.0237546805,\qquad
  \theta=2.1426\pm 2.8795\,i,
\end{equation}
with eigenvector $(0.3557\pm 0.3776\,i,\;0.8549)$, obtained from the numerical
Jacobian of \eqref{eq:solveeta}--\eqref{eq:betaxi0}; $\eta=2$ at the fixed
point.  The paper's own closed-form stability-matrix entries $M_{00}$ do not
equal the numerical derivatives of the appendix flows at the GMFP
($M_{\lambda_0\lambda_0}$ differs by exactly $+4$, $M_{\xi_0\xi_0}$ by
$\approx-2.14$); the published exponents and eigenvectors are reproduced by
the flow-derived Jacobian, and the paper itself notes transcription errors in
its precursor.  We therefore use the closed-form flows and numerical
Jacobians only, and record the discrepancy.

The scaffold wiring evaluates the $\xi_2$ flow per species on the GMFP
background of \eqref{eq:gmfp} with its exact fixed-point value $\eta=2$, and
maps
\begin{equation}
  \label{eq:npwiring}
  \lambda_2=m_i^2,\qquad \xi_2=\xi_i,\qquad
  \lambda_4=\lambda_i/4
  \qquad(i=H,\Phi,E),
\end{equation}
from $V=\lambda_i/4\,(\phi^2)^2$ for the radial field.  Because the
background is fixed at the NP GMFP, the $\xi$ entries do not mix the NP gauge
with the DEP gravity block; the two convention labels never coexist in one
solve.  At Gaussian matter ($\lambda_2=\xi_2=\lfour=0$) the $\xi$ flow
vanishes (residual $\sim 10^{-16}$), so $\xi_i^*=0$ is an exact fixed point of
this block.

\subsection{The scalar local-potential block: masses, quartics, wave
functions}
\label{sec:lpa}

The six real scalars (four Higgs components, $\Phi_c$, $E$) are described by a
local potential $U_k(h,\Phi,E)$ whose Hessian on a constant background
$\varphi$ is
\begin{equation}
  \label{eq:hessian}
  m^2_{ab}(\varphi)=m^2_{ab}(0)
  +\tfrac12\,\lambda_{abcd}\,\varphi_c\varphi_d
  +\tfrac{1}{24}\,c_{abcdef}\,\varphi_c\varphi_d\varphi_e\varphi_f ,
\end{equation}
where $\lambda_{abcd}$ is the quartic tensor and $c_{abcdef}$ the sextic
tensor of App.~\ref{app:sextic}.  With the optimized regulator
\eqref{eq:litim}, the scalar contribution to the flow of the potential is the
constant-background trace
\begin{equation}
  \label{eq:trace}
  \dt U_k\big|_{\mathrm{scalar}}
  =\frac{k^4}{32\pi^2}\sum_{a=1}^{6}\Big(1+\frac{m^2_a(\varphi)}{k^2}\Big)^{-1},
\end{equation}
with $m^2_a$ the eigenvalues of \eqref{eq:hessian}.  Projecting at zero field
gives the quantum contributions to the masses and quartics,
\begin{equation}
  \label{eq:quarticprojections}
  \beta_{m_i^2}\supset \frac{\partial^2 \dt U_k}{\partial \varphi_i^2}\Big|_0,\qquad
  \beta_{\lambda_i}\supset \frac{1}{6}\,\frac{\partial^4 \dt U_k}
  {\partial \varphi_i^4}\Big|_0,\qquad
  \beta_{\kappa_{ij}}\supset \frac{\partial^4 \dt U_k}
  {\partial \varphi_i^2\,\partial\varphi_j^2}\Big|_0 ,
\end{equation}
computed with mixed finite-difference stencils whose step is checked against
an analytic toy derivative.  The Gaussian point is an exact fixed point of
this block, and the fourth-order projections match the analytically known
one-scalar flow.  The wave-function coordinates are driven by the DEP
anomalous dimensions of Sec.~\ref{sec:dep},
\begin{equation}
  \label{eq:Zflow}
  \beta_{Z_i}=-\eta_s\,Z_i \qquad (i=h,\phi,E),
\end{equation}
with $\eta_s\simeq -0.076$ at the Standard-Model fixed point, so $Z_i^*=0$.

\subsection{Dimension-six scalar potentials (sextic LPA)}
\label{sec:sextic}

The ten scalar dimension-six potentials are projected from the same trace
\eqref{eq:trace} with the sixth-order potential.  The normalization
convention assigns a factor $1/2^{a+b+c}$ to every monomial
$(h^2)^a(\Phi^2)^b(E^2)^c$ with $a+b+c=3$, the same ``two per invariant
factor'' rule that fixes the quartic sector ($\lambda_H/4\,(h^2)^2$,
$\kappa_{H\Phi}/4\,(h^2)\Phi^2$).  The projections are
\begin{equation}
  \label{eq:sexticprojections}
  c_{H6}=\frac{1}{90}\,\frac{\partial^6\dt U_k}{\partial h^6}\Big|_0,\qquad
  c_{H4\Phi2}=\frac{1}{6}\,\frac{\partial^4\partial^2\dt U_k}
  {\partial h^4\,\partial\Phi^2}\Big|_0,\qquad
  c_{H2\Phi2E2}=\frac{\partial^2\partial^2\partial^2\dt U_k}
  {\partial h^2\,\partial\Phi^2\,\partial e^2}\Big|_0,
\end{equation}
and cyclically; the $O(4)$ combinatorics of the Higgs sector are fixed by the
sextic tensor of App.~\ref{app:sextic}, which is validated by a unit test
that reconstructs every monomial of the potential from the tensor at
generic couplings.  The finite-difference step is chosen in the convergence
plateau ($\delta\simeq 5\times10^{-2}$; values at $\delta=3\times10^{-2}$ and
$7\times10^{-2}$ agree to a few percent).  The block passes three controls:
it reduces exactly to the quartic flow at vanishing sextics; the pure-toy
derivatives $\partial^6/\partial h^6=720$, $\partial^4\partial^2=144$,
$\partial^2\partial^2\partial^2=48$ are reproduced; and all ten entries
vanish at the Gaussian point to the stencil noise floor
($\sim 10^{-11}$).

\subsection{The curvature-squared sector and its negative result}
\label{sec:r2}

The $R^2$ sector follows Codello, Percacci and Rahmede
\cite{cpr09} (arXiv:0805.2909), one-loop approximation with the type-III
generalized optimized cutoff and the Shapiro gauge choice that makes the
quadratic action minimal fourth order.  The action ansatz is
\begin{equation}
  \label{eq:cpr}
  S=\int \sqrt{g}\,\Big[
    \frac{1}{2\lambda}\,C^2+\frac{1}{\xi}\,R^2-\frac{1}{\rho}\,E\Big],
\end{equation}
with $C^2$ the Weyl-squared, $E$ the Euler density, and beta functions
\begin{equation}
  \label{eq:cprbetas}
  \beta_\lambda=-\frac{133}{10}\frac{\lambda^2}{(4\pi)^2},\qquad
  \beta_\xi=-\frac{10\lambda^2-5\lambda\xi+5\xi^2/36}{(4\pi)^2},\qquad
  \beta_\rho=\frac{196}{45}\frac{\rho^2\lambda}{(4\pi)^2}.
\end{equation}
In the ratios $\omega=-3\lambda/\xi$ and $\theta=\lambda/\rho$ the fixed
points are $\omega_*\in\{-5.467,-0.0228\}$ and $\theta_*=56/171\simeq 0.3275$;
the $\omega$ conversion is exactly consistent with \eqref{eq:cprbetas} via the
chain rule (checked algebraically and numerically).  We record the
paper-internal caveat that the printed $\beta_\theta$ does not follow from
$\beta_\rho$ by the chain rule; the $\theta$ sector is decoupled at this
order, so the discrepancy affects no wired coordinate, and
$\theta_*=56/171$ is used as printed.

The scaffold coordinates are the inverse couplings
\begin{equation}
  \label{eq:scaffoldr2}
  a_{R^2}=\frac{1}{\xi},\qquad a_{C^2}=\frac{1}{2\lambda},
\end{equation}
defined only for $\lambda\neq 0$ (singular at the Gaussian point).  From
\eqref{eq:cprbetas},
\begin{equation}
  \label{eq:scaffoldbetas}
  \beta_{a_{C^2}}=\frac{133}{320\,\pi^2}\approx 0.04211\;(>0,\ \text{constant}),
  \qquad
  \beta_{a_{R^2}}=\frac{10x^2-5x+5/36}{(4\pi)^2},\qquad
  x=\frac{a_{R^2}}{2a_{C^2}}.
\end{equation}
$\beta_{a_{R^2}}$ vanishes on the two lines $a_{R^2}=2x\,a_{C^2}$ with
$x\in\{0.0295207,\,0.4704793\}$, but $\beta_{a_{C^2}}$ is strictly positive
and constant, so the Weyl-squared direction runs monotonically and the
$(a_{R^2},a_{C^2})$ subsystem has no fixed point in this parametrization.  We
record this as a negative result, not a papered-over gap.  Physically, the
underlying flow \eqref{eq:cprbetas} has its Gaussian point at
$\lambda=\xi=0$ (the Weyl-squared coupling is asymptotically free in the
one-loop approximation), and the inverse parametrization diverges there; the
conclusion is a statement about the scaffold coordinates, which are the ones
used in the $41$-coordinate basis.

\section{The joint fixed point}
\label{sec:joint}

\subsection{Fixed-point vector and residuals}
\label{sec:fpvector}

Assembling the blocks of Sec.~\ref{sec:blocks} into
\texttt{beta\_partial} (33 of the 41 entries finite; the remaining 8 are
\emph{pending}, Sec.~\ref{sec:coverage}), the joint fixed point of the
derived flow is
\begin{equation}
  \label{eq:jointfp}
  (\lambda_N,\,g_N)=(-4.7975095,\,1.7621271),
  \qquad
  \text{all other 31 derived coordinates }=0,
\end{equation}
i.e.\ the DEP Standard-Model row \eqref{eq:depfp} with the matter sector at
its Gaussian point.  The residuals at \eqref{eq:jointfp} are: gravity and
matter entries zero to machine precision; $\xi_i$ entries $\sim 10^{-16}$
(exact fixed point of the NP block, Sec.~\ref{sec:np}); dimension-six entries
$\sim 10^{-11}$ (the finite-difference noise floor of the sextic
projections).  The joint fixed point is therefore a genuine fixed point of
the derived flow to the numerical accuracy of the projections.

\subsection{The 33-dimensional stability spectrum}
\label{sec:spectrum}

The stability matrix is the numerical central-difference Jacobian of the
derived flow restricted to the 33 derived coordinates, evaluated at
\eqref{eq:jointfp}.  Two conventions are reported: the quantum-only
Jacobian of \texttt{beta\_partial}, and the quantum-plus-canonical Jacobian
with the canonical terms \eqref{eq:canonical} added for the mass-squared and
dimension-six entries (the DEP gravity block contains its own canonical
terms; all other matter couplings are dimensionless).  The
quantum-plus-canonical spectrum is the physically relevant one and is
displayed in Table~\ref{tab:spectrum}.

\begin{table}[htbp]
  \centering
  \caption{Quantum-plus-canonical critical exponents
  $\theta=-\mathrm{eig}(\partial\beta/\partial g)$ of the 33-dimensional
  derived flow at the joint fixed point \eqref{eq:jointfp}.}
  \label{tab:spectrum}
  \begin{tabular}{lcc l}
    \toprule
    $\theta$ & multiplicity & class & directions \\
    \midrule
    $+3.9609$   & 1 & relevant & $\Lambda$-direction ($\lambda_N$) \\
    $+2.8554$   & 3 & relevant & $m_H^2$, $m_\Phi^2$, $m_E^2$ \\
    $+2.0000$   & 3 & relevant & $\xi_H$, $\xi_\Phi$, $\xi_E$ \\
    $+1.6436$   & 1 & relevant & $\tilde G$-direction ($g_N$) \\
    $0.0000$    & 12 & marginal & $g_1,g_2,g_3,y_t,y_b,y_\tau$ and the six quartics \\
    $-0.0756$   & 3 & irrelevant & $Z_h,Z_\phi,Z_E$ \\
    $-2.0000$   & 10 & irrelevant & the ten dimension-six potentials \\
    \bottomrule
  \end{tabular}
\end{table}

The two gravity eigenvalues reproduce the DEP table-convention exponents
$(3.961,1.644)$ of \eqref{eq:depfp} exactly.  The counting is
\begin{equation}
  \label{eq:counting}
  \underbrace{8}_{\text{relevant}}\;+\;
  \underbrace{12}_{\text{marginal}}\;+\;
  \underbrace{13}_{\text{irrelevant}}\;=\;33 .
\end{equation}

\subsection{Interpretation}
\label{sec:interpretation}

The spectrum has a transparent block structure.
\begin{itemize}
  \item \emph{Gravity.} The two exponents $+3.9609$ and $+1.6436$ are the
        DEP eigenvalues, unchanged by the (Gaussian) matter sector; they
        confirm that the fixed point is UV attractive in the
        Einstein--Hilbert directions.
  \item \emph{Mass-squared directions.} $\theta=2.8554=2+0.8554$: the
        canonical weight $+2$ (from $-[m_i^2]=+2$ per
        \eqref{eq:canonical}, i.e.\ $\beta_{m^2}\supset -2m^2$) is augmented
        by a quantum shift from the scalar one-loop traces of
        Secs.~\ref{sec:lpa}--\ref{sec:sextic} evaluated at the Standard-Model
        row of the DEP block.  The three masses are relevant with identical
        exponents.
  \item \emph{Nonminimal couplings.} $\theta=2.0000$ exactly for all three
        $\xi_i$: at the NP GMFP background with $\eta=2$ the $\xi_2$ flow
        \eqref{eq:xi2implicit} has an exact linear coefficient, and the
        $\xi_i$ directions inherit it.  This is a pure quantum result; the
        $\xi_i$ have canonical weight zero.
  \item \emph{Marginal matter couplings.} The twelve coordinates
        $g_1,g_2,g_3,y_t,y_b,y_\tau$ and the six quartic couplings
        $\lambda_H,\lambda_\Phi,\lambda_E,\kappa_{H\Phi},\kappa_{HE},
        \kappa_{\Phi E}$ have vanishing eigenvalues: at Gaussian matter every
        one-loop contribution is of higher order in the couplings (the
        quartic and Yukawa beta functions have no linear terms in the
        respective couplings, and the gauge beta functions vanish at zero
        gauge coupling), and the canonical weights are zero.  The stability
        matrix does not lift their degeneracy at this point, a well-known
        feature of Gaussian matter sectors.
  \item \emph{Wave functions.} $\theta=-0.0756=-\eta_s$ with
        $\eta_s\simeq-0.076$ from \eqref{eq:Zflow}; the $Z_i$ are irrelevant
        and the fixed point is $Z_i^*=0$.
  \item \emph{Dimension-six potentials.} $\theta=-2.0000$ exactly: pure
        canonical scaling \eqref{eq:canonical}; the quantum parts vanish at
        Gaussian matter to the stencil floor, so all ten operators are
        irrelevant at this fixed point.
\end{itemize}

The structure is consistent with the expectation for an interacting
gravity--matter fixed point with Gaussian matter: the gravity sector
controls the flow, the masses and nonminimal couplings are relevant, and all
higher-dimensional scalar operators are irrelevant.

\section{Coverage, limitations, and pending directions}
\label{sec:coverage}

Thirty-three of the 41 coordinates are derived and verified:
\begin{itemize}
  \item \texttt{beta\_partial} (quantum-only, DEP gravity + SM matter +
        scalar LPA + NP $\xi$ + sextic LPA): entries 0--1 (gravity),
        5--19 (SM gauge/Yukawa/quartics/portals), 20--22 ($\xi_i$),
        23--25 ($Z_i$), 26--35 (dimension-six potentials), plus the scalar
        LPA mass entries 11,13,15.
  \item The remaining coordinates are \emph{pending} and return \texttt{NaN}:
        $a_{R^2},a_{C^2}$ (handled by the negative result of
        Sec.~\ref{sec:r2}); the derivative operators $c_{H\Box},c_{\Phi\Box},
        c_{E\Box}$ (not projectable by a constant-background LPA); and the
        gauge-strength portals $c_{HF^2},c_{\Phi F^2},c_{EF^2}$ (require a
        dressed-gauge-sector trace).
\end{itemize}
No flow value is fabricated for the pending coordinates, and the framework
refuses to present a $41$-dimensional fixed-point claim.  Every numerical
result is accompanied by its convention label, and the unit-test suite pins
each verified number.

\section{Discussion}
\label{sec:discussion}

The reconstruction illustrates a methodological point that is often
understated: in the FRG, a coupling is not ``there'' until its projected
flow is written down.  The $41$-coordinate claim that motivated this project
contains no operator list, no beta functions, no fixed-point vector, and no
stability matrix; we have therefore delivered a smaller system (33 derived
coordinates) in which every object is written down, sourced, and verified.
The main physics findings are:

\begin{enumerate}
  \item The joint fixed point of the derived flow is the DEP
        Standard-Model fixed point with Gaussian matter.  Its spectrum has
        exactly 8 relevant, 12 marginal, and 13 irrelevant directions.  The
        gravity eigenvalues are the DEP values $(3.961,1.644)$; the masses
        and nonminimal couplings are relevant; the ten dimension-six scalar
        operators are irrelevant by canonical scaling alone.
  \item The $\xi_i$ directions are exact: the Narain--Percacci GMFP
        background with $\eta=2$ forces $\theta=2$ and an exact Gaussian
        fixed point of the nonminimal sector.  This is a clean, verifiable
        sub-result and a nontrivial consistency check between the two
        published blocks.
  \item The curvature-squared sector has no fixed point in the scaffold
        parametrization \eqref{eq:scaffoldr2}: $\beta_{a_{C^2}}>0$ constant.
        In the original couplings the Weyl-squared direction is asymptotically
        free toward the Gaussian point; the negative result is a property of
        the inverse coordinates used in the basis.  We record it explicitly,
        because an unguarded reader of the scaffold might otherwise search
        for a fixed point that does not exist.
  \item The twelve marginal matter couplings (gauge, Yukawa, quartic) at
        Gaussian matter are not a prediction of fixed values; they are the
        statement that the stability matrix is degenerate in those
        directions at the Gaussian point.  A fully interacting matter fixed
        point would require the coupled matter beta functions away from the
        Gaussian point, which is outside the currently derived system.
\end{enumerate}

\emph{Limitations.} (i) All gravity and matter blocks are one-loop
(threshold-functional) or LPA-level; there is no fermion trace in the scalar
projections, no dressed gauge sector, and no derivative (curvature-squared or
box-type) projection for matter.  (ii) The pending coordinates are not
fabricated; the joint fixed point is a statement about the 33 derived
coordinates only.  (iii) The conventions (type-II vs.\ optimized regulator,
de-Donder vs.\ type-A gauge, table vs.\ improved ghost factors) are
deliberately separated; mixing them would produce numbers that are neither
paper's.  (iv) The NP $M_{00}$ discrepancy and the CPR $\beta_\theta$
inconsistency are recorded paper-internal observations; both are resolved in
favor of the closed-form flows and decoupled sectors.

\section{Conclusions}
\label{sec:conclusions}

Within a fully declared truncation, the reconstructed gravity--matter system
has a UV-attractive joint fixed point with 8 relevant, 12 marginal, and 13
irrelevant directions at the DEP Standard-Model row with Gaussian matter.  All
33 derived flows are traced to published sources or documented projections and
are pinned by unit tests; the 8 remaining coordinates are explicitly pending,
and the curvature-squared sector is documented as having no fixed point in the
scaffold parametrization.  The accompanying implementation
(\texttt{outputs/uv\_reconstruction\_41/}) is the precise definition of the
calculation; this paper is the human-readable record of it.

\section*{Reproducibility}
\label{sec:repro}
The complete implementation lives in
\texttt{outputs/uv\_reconstruction\_41/}.  The verified numbers of this paper
are regenerated by
\begin{verbatim}
  $PY -m unittest discover -s tests
\end{verbatim}
(44 tests), with \texttt{PY = /Users/christophermichaelbaird/.cache/}\\[-2pt]
\texttt{codex-runtimes/codex-primary-runtime/dependencies/python/bin/python3}.
The joint fixed point and spectrum are produced by
\texttt{joint\_fixed\_point\_stability.py}; the convention register is
\texttt{DECLARED\_TRUNCATION\_V0\_1.md}.

\begin{thebibliography}{99}

\bibitem{wetterich} C.~Wetterich, \emph{Exact evolution equation for the
effective potential}, Phys.\ Lett.\ B 301 (1993) 90, arXiv:hep-ph/9211263.

\bibitem{berges} J.~Berges, N.~Tetu, C.~Wetterich,
\emph{The effective average action and its flow equation}, Phys.\ Rept.\ 363
(2002) 223, arXiv:hep-ph/0005122.

\bibitem{reuter} M.~Reuter, \emph{Nonperturbative evolution equation for
quantum gravity}, Phys.\ Rev.\ D 57 (1998) 971, arXiv:hep-th/9609127.

\bibitem{litim01} D.~F.~Litim, \emph{Optimised renormalisation group flows},
Phys.\ Rev.\ D 64 (2001) 105007, arXiv:hep-th/0009045.

\bibitem{litim06} D.~F.~Litim, \emph{Fixed points of quantum gravity and the
renormalisation group}, Phys.\ Rev.\ D 76 (2007) 125001,
arXiv:hep-th/0606044.

\bibitem{dep14} P.~Don\`a, A.~Eichhorn, R.~Percacci, \emph{Matter matters in
asymptotically safe quantum gravity}, Phys.\ Rev.\ D 89 (2014) 084035,
arXiv:1311.2898.

\bibitem{narain10} G.~Narain, R.~Percacci, \emph{On the stability of the
asymptotic solutions of the Asymptotic Safety program}, Class.\ Quantum
Grav.\ 27 (2010) 075001, arXiv:0911.0386.

\bibitem{cpr09} A.~Codello, R.~Percacci, C.~Rahmede, \emph{Investigating the
ultraviolet properties of gravity with a Wilsonian renormalization group
equation}, Ann.\ Phys.\ (NY) 324 (2009) 414, arXiv:0805.2909.

\bibitem{percacciBook} R.~Percacci, \emph{An Introduction to Covariant
Quantum Gravity and Asymptotic Safety}, World Scientific (2017).

\end{thebibliography}

\appendix

\section{The declared 41-coordinate operator basis}
\label{app:operators}

\begin{table}[htbp]
  \centering
  \caption{The declared 41-coordinate basis.  $[g_i]$ is the mass dimension
  of the coupling; the canonical term is $-[g_i]\,g_i$.}
  \label{tab:operators}
  \small
  \begin{tabular}{llc l}
    \toprule
    \# & coordinate & $[g_i]$ & operator \\
    \midrule
    1  & $g_N$       & $-2$ & $R$ (dimensionless Newton coupling) \\
    2  & $\lambda_N$ & $+2$ & cosmological term \\
    3  & $a_{R^2}$   & $-$ & $R^2$ \\
    4  & $a_{C^2}$   & $-$ & Weyl-squared \\
    5--7  & $g_1,g_2,g_3$ & $0$ & SM gauge kinetic terms \\
    8--10 & $y_t,y_b,y_\tau$ & $0$ & third-generation Yukawas \\
    11,13,15 & $m_H^2,m_\Phi^2,m_E^2$ & $+2$ & scalar masses \\
    12,14,16 & $\lambda_H,\lambda_\Phi,\lambda_E$ & $0$ & scalar quartics \\
    17 & $\kappa_{\Phi E}$ & $0$ & $\Phi_c^2 E^2$ \\
    18--19 & $\kappa_{H\Phi},\kappa_{HE}$ & $0$ & $|H|^2\Phi_c^2$, $|H|^2E^2$ \\
    20--22 & $\xi_H,\xi_\Phi,\xi_E$ & $0$ & $R|H|^2$, $R\Phi_c^2$, $RE^2$ \\
    23--25 & $Z_h,Z_\phi,Z_E$ & $0$ & wave-function renormalizations \\
    26 & $c_{H6}$ & $-2$ & $|H|^6$ \\
    27 & $c_{\Phi6}$ & $-2$ & $\Phi^6$ \\
    28 & $c_{E6}$ & $-2$ & $E^6$ \\
    29--34 & $c_{H4\Phi2},\ldots,c_{\Phi2E4}$ & $-2$ & mixed sixth-order potentials \\
    35 & $c_{H2\Phi2E2}$ & $-2$ & $|H|^2\Phi^2E^2$ \\
    36--38 & $c_{H\Box},c_{\Phi\Box},c_{E\Box}$ & $-2$ & $|H|^2\Box|H|^2$, $\Phi^2\Box\Phi^2$, $E^2\Box E^2$ (pending) \\
    39--41 & $c_{HF^2},c_{\Phi F^2},c_{EF^2}$ & $-2$ & $|H|^2F^2$, $\Phi^2F^2$, $E^2F^2$ (pending) \\
    \bottomrule
  \end{tabular}
\end{table}

\section{The Narain--Percacci five-coupling system}
\label{app:np}

With $\lambda_0,\xi_0,\lambda_2,\xi_2,\lfour$ as in Sec.~\ref{sec:np},
$\eta$ from \eqref{eq:solveeta}, $D\equiv\xi_0-\lambda_0$,
$A\equiv 1+2\lambda_2$, $B\equiv 1+2\xi_2$, and
$\dot\xi_2\equiv\beta_{\xi_2}$, the transcribed closed forms are
\begin{align}
  Z&=\frac{1}{576\pi^2}\Bigg[
    \frac{A}{D}\Big(9+\frac{39\xi_0}{D}+\frac{60\xi_0^2}{D^2}\Big)
    +\frac{3(3+32\xi_2)}{D}
    -\frac{6\xi_0(11+2\xi_2)}{D^2}
    -\frac{60\xi_0^2 B}{D^3}
    +\frac{216\xi_2 B^2}{A^3 D}\nonumber\\
  &\qquad\quad+\frac{9\big(\lambda_0(5-2\xi_2)-2\xi_0B\big)B}{A D^2}
    +\frac{27B(1-10\xi_2-16\xi_2^2)}{A^2 D}
    +\frac{108\xi_0\xi_2B^2}{A^2 D^2}\nonumber\\
  &\qquad\quad+\frac{72\lfour(1+12\xi_2+2\lambda_2)}{A^3}\Bigg],\\
  W&=\frac{1}{1152\pi^2}\Bigg[
    \frac{A}{D}\Big(3+\frac{18\xi_0}{D}+\frac{20\xi_0^2}{D^2}\Big)
    +\frac{15\xi_2}{\xi_0}-\frac{6(1+\xi_2)}{D}
    -\frac{10\xi_0(3+4\xi_2)}{D^2}
    -\frac{20\xi_0^2B}{D^3}\nonumber\\
    &\qquad-\frac{3\big(\lambda_0-\xi_0(5-4\xi_2)\big)B}{A D^2}
    +\frac{36\xi_0\xi_2B^2}{A^2 D}\Bigg],\\
  V&=\frac{1}{1152\pi^2}\Bigg[
    -15+\frac{54\xi_0}{D}+\frac{20\xi_0^2}{D^2}
    -\frac{6\xi_0(7+2\xi_2)}{A D}\Bigg]
    \nonumber\\
  &\quad-\frac{1}{1152\pi^2}\cdot\frac{144\,\xi_0\xi_2\,B}{A^2 D},
\end{align}
with $\beta_{\xi_2}$ the solved implicit form
\begin{equation}
  \label{eq:xi2solved}
  \beta_{\xi_2}=\frac{Z+\eta W}{1-V/\xi_0},
\end{equation}
and the remaining flows
\begin{align}
  \beta_{\lambda_2}&=-2\lambda_2+\frac{1}{48\pi^2}\Bigg[
    \frac{9\lambda_0 B}{2D^2}-\frac{9(2\lambda_0-\xi_0)B^2}{2 A D^2}\Bigg]
    \nonumber\\
  &\qquad-\frac{1}{48\pi^2}\Bigg[
    \frac{9B^2}{2 A^2 D}+\frac{18\lfour}{A^2}\Bigg]\nonumber\\
    &\quad+\frac{\eta}{96\pi^2}\Bigg[-\frac{2\xi_2}{\xi_0}
    +\frac{3\xi_0 B}{2D^2}-\frac{3\xi_0 B^2}{2 A D^2}\Bigg]
    +\frac{\dot\xi_2}{\xi_0}\cdot\frac{1}{96\pi^2}\Bigg[2
    -\frac{3\xi_0}{D}+\frac{6\xi_0 B}{A D}\Bigg],\\
  \beta_{\lfour}&=\frac{1}{48\pi^2}\Bigg[
    \frac{9}{4D^2}\Big(5A(1+4\xi_2)-B(21+62\xi_2)
    +\frac{33B^3}{A}-\frac{B^3(23+24\xi_2)}{A^2}+\frac{6B^4}{A^3}\Big)\nonumber\\
    &\quad+\frac{9\xi_0(\xi_2-\lambda_2)^2}{D^3}\,q
    -\frac{72\lambda_2\lfour B(1-4\lambda_2+6\xi_2)}{D A^3}
    +\frac{9\xi_0\lfour\,q}{D^2}+\frac{216\lfour^2}{A^3}\Bigg]\nonumber\\
    &\quad+\frac{\eta}{96\pi^2}\Bigg[\frac{2\xi_2^2}{\xi_0^2}
    +\frac{3\xi_0(\xi_2-\lambda_2)^2}{D^3}\,q
    +\frac{3\xi_0\lfour\,q}{D^2}\Bigg]\nonumber\\
    &\quad+\frac{\dot\xi_2}{\xi_0}\cdot\frac{1}{96\pi^2}\Bigg[
    -\frac{2\xi_2}{\xi_0}-\frac{3\xi_0(\xi_2-\lambda_2)}{D^2}
    \Big(12\frac{B^2}{A^2}-21\frac{B}{A}+10\Big)
    -\frac{24\xi_0\lfour(1-4\lambda_2+6\xi_2)}{A^2 D}\Bigg],
\end{align}
where
\begin{equation}
  q=6\frac{B^2}{A^2}-10\frac{B}{A}+5 .
\end{equation}
These equations reproduce the paper's GMFP, exponents, and eigenvector
(Sec.~\ref{sec:np}).

\section{The sextic tensor and its projections}
\label{app:sextic}

With the potential written $V=c_{abcdef}\,\phi^a\phi^b\phi^c\phi^d\phi^e
\phi^f/6!$ and the monomial convention of Sec.~\ref{sec:sextic}, the
$O(4)\times O(1)\times O(1)$ sextic tensor is generated from the ten
coordinates by the following symmetric assignments (indices $i,j,k\in\{1..4\}$
label Higgs components, $\Phi=5$, $E=6$; each line is assigned to every
distinct index pattern):
\begin{align}
  \label{eq:tensor}
  c_{iiiiii}&=90\,c_{H6},\qquad c_{\Phi\Phi\Phi\Phi\Phi\Phi}=90\,c_{\Phi6},
  \qquad c_{EEEEEE}=90\,c_{E6},\\
  c_{iiiijj}&=18\,c_{H6},\qquad c_{iijjkk}=6\,c_{H6},\\
  c_{iiii\Phi\Phi}&=6\,c_{H4\Phi2},\qquad c_{iijj\Phi\Phi}=2\,c_{H4\Phi2},
  \qquad c_{iiiiEE}=6\,c_{H4E2},\qquad c_{iijjEE}=2\,c_{H4E2},\\
  c_{ii\Phi\Phi\Phi\Phi}&=6\,c_{H2\Phi4},\qquad c_{iiEEEE}=6\,c_{H2E4},\\
  c_{\Phi\Phi\Phi\Phi EE}&=6\,c_{\Phi4E2},\qquad c_{\Phi\Phi EEEE}=6\,c_{\Phi2E4},\\
  c_{ii\Phi\Phi EE}&=c_{H2\Phi2E2}.
\end{align}
A unit test reconstructs every monomial from the tensor at generic couplings,
guaranteeing the convention is enforced.  The projections of the trace
\eqref{eq:trace} onto the ten operators are
\begin{align}
  c_{H6}&=\frac{1}{90}\,\frac{\partial^6}{\partial h^6}\,\dt U_k\Big|_0,
  \qquad c_{\Phi6}=\frac{1}{90}\,\frac{\partial^6}{\partial\Phi^6}\,\dt U_k\Big|_0,
  \qquad c_{E6}=\frac{1}{90}\,\frac{\partial^6}{\partial e^6}\,\dt U_k\Big|_0,\\
  c_{H4\Phi2}&=\frac{1}{6}\,\frac{\partial^6}{\partial h^4\partial\Phi^2}
    \dt U_k\Big|_0,\qquad\text{etc.},\qquad
  c_{H2\Phi2E2}=\frac{\partial^6}{\partial h^2\partial\Phi^2\partial e^2}
    \dt U_k\Big|_0,
\end{align}
computed with mixed finite-difference stencils at step
$\delta=5\times10^{-2}$ in the convergence plateau.

\end{document}
